\chapter{Diseases of the Breast}
\label{ch:breast}
The breast is a modified sweat gland consisting of glandular tissue (lobules and ducts), fibrous connective tissue, and adipose tissue. It is present in both sexes but is functionally developed in females under hormonal influence. Breast diseases include benign conditions (fibrocystic changes, fibroadenomas, mastitis), pre-malignant conditions, and breast cancer---the most common cancer and second leading cause of cancer-related death in women worldwide.

%-------------------------------------------------------------
\section{Benign Breast Conditions}
%-------------------------------------------------------------

%-------------------------------------------------------------

%-------------------------------------------------------------

%-------------------------------------------------------------

%-------------------------------------------------------------

%-------------------------------------------------------------

%-------------------------------------------------------------

%-------------------------------------------------------------

\label{sec:Benign Breast Conditions}
%-------------------------------------------------------------%-------------------------------------------------------------

\subsection{Fibroadenoma}
\label{sec:Fibroadenoma}
Fibroadenomas are the most common benign breast tumors in young women (ages 15--35). The condition results from estrogen-dependent benign neoplasms composed of fibrous and glandular tissue. Clinical features typically include a solitary, firm, mobile, well-circumscribed, rubbery mass---often described as a ``breast mouse'' due to its mobility---which can grow during pregnancy or hormonal stimulation. Diagnosis is by ultrasound (well-circumscribed, hypoechoic, oval mass with horizontal orientation) and core needle biopsy to confirm. Management includes surgical excision for giant fibroadenomas or those causing symptoms; small, asymptomatic fibroadenomas may be observed. Prognosis is excellent.

\subsection{Fibrocystic Changes}
\label{sec:Fibrocystic Changes}
Fibrocystic changes are the most common benign breast condition, affecting up to 60\% of women at some point. The condition is characterized by breast pain (mastalgia), tenderness, nodularity, and thickening, often fluctuating with the menstrual cycle, and is hormone-dependent, typically worsening in the luteal phase. Histologically, it encompasses a spectrum of findings: cysts (microcystic and macrocystic), fibrosis, adenosis, and epithelial hyperplasia. Clinical features include breast pain and nodularity. Diagnosis is by clinical breast examination, mammography (for women over 40), and ultrasound (to evaluate dominant cysts). Management includes reassurance, well-fitting support bras, NSAIDs, and for severe symptoms, oral contraceptives or tamoxifen. Prognosis is excellent, with no malignant potential.

\subsection{Gynecomastia}
\label{sec:Gynecomastia}
Gynecomastia is benign enlargement of breast tissue in males, caused by an imbalance of estrogen and androgen effects, present in up to 70\% of adolescent males (physiologic gynecomastia, typically resolving within 1--2 years) and up to 65\% of men over 50. Pathological causes include medications (spironolactone, finasteride, cimetidine, chemotherapy, exogenous testosterone/steroids, cannabinoids), liver disease, hypogonadism, hyperthyroidism, Klinefelter syndrome, and obesity (pseudogynecomastia). Clinical features include a palpable, firm, disk-shaped tissue beneath the areola. Diagnosis is by clinical examination and evaluation of hormonal levels. Management includes addressing reversible causes, observation for physiologic or medication-induced gynecomastia, and surgical mastectomy for persistent, symptomatic cases. Prognosis is excellent.

\subsection{Mastitis and Breast Abscess}
\label{sec:Mastitis and Breast Abscess}
Mastitis is inflammation of the breast tissue, most commonly lactational (puerperal) mastitis, occurring in 2--10\% of breastfeeding women. The condition results from bacterial infection, most commonly \textit{Staphylococcus aureus}, introduced through a cracked nipple or cracked skin. Clinical features include a focal area of erythema, warmth, tenderness, and induration, often accompanied by fever and malaise. Diagnosis is clinical. Management includes continued breastfeeding or expression, warm compresses, and antibiotics (dicloxacillin or cephalexin). Breast abscess is a complication of mastitis requiring aspiration or incision and drainage. Prognosis is excellent with appropriate treatment.

\subsection{Phyllodes Tumor}
\label{sec:Phyllodes Tumor}
Phyllodes tumors are rare fibroepithelial neoplasms that can be benign, borderline, or malignant (10--15\% of all phyllodes tumors), typically presenting as large, rapidly growing, painless breast masses, more common in women 40--50 years of age. The condition results from stromal proliferation in fibroepithelial tissue. Histologically, they are distinguished from fibroadenomas by increased stromal cellularity, stromal overgrowth, atypia, and mitotic activity. Clinical features include a rapidly enlarging breast mass. Diagnosis is by core needle biopsy and definitive surgical excision with wide margins ($>$1 cm) is required for all phyllodes tumors, regardless of grade. Prognosis is good with complete excision; malignant phyllodes tumors can metastasize hematogenously.

%-------------------------------------------------------------
\section{Breast Cancer}
%-------------------------------------------------------------

%-------------------------------------------------------------

%-------------------------------------------------------------

%-------------------------------------------------------------

%-------------------------------------------------------------

%-------------------------------------------------------------

%-------------------------------------------------------------

%-------------------------------------------------------------

\label{sec:section:Breast Cancer}
%-------------------------------------------------------------%-------------------------------------------------------------

\subsection{Breast Cancer}
\label{sec:Breast Cancer}
Breast cancer is the most commonly diagnosed cancer worldwide (2.3 million new cases annually) and the second leading cause of cancer-related death in women. Risk factors include female sex, advancing age, family history (particularly BRCA1 and BRCA2 mutations---lifetime risk 45--72\% and 26--43\% respectively), prior breast cancer, early menarche, late menopause, nulliparity, late first pregnancy, obesity (postmenopausal), alcohol use, and prior chest radiation, with protective factors including breastfeeding, physical activity, and early menopause. The majority of breast cancers (70--80\%) are invasive ductal carcinoma (not otherwise specified), with other types including invasive lobular carcinoma (10--15\%), inflammatory breast cancer, Paget's disease of the nipple, and male breast cancer ($<$1\%). The condition results from malignant transformation of breast epithelial cells.

Screening: mammography is recommended biennially for women ages 50--74 (USPSTF) or annually starting at age 40 (American College of Radiology), with earlier and more intensive screening for high-risk individuals including breast MRI.

Clinical features include a painless breast mass, skin changes, nipple discharge, and axillary lymphadenopathy. Diagnosis: diagnostic mammography with ultrasound; core needle biopsy for histopathological evaluation. Molecular subtyping by immunohistochemistry: estrogen receptor (ER), progesterone receptor (PR), HER2, and Ki-67, with four molecular subtypes: luminal A (ER+/PR+, HER2-, low Ki-67---best prognosis), luminal B (ER+, HER2+ or high Ki-67), HER2-enriched (ER-, PR-, HER2+), and triple-negative/basal-like (ER-, PR-, HER2----worst prognosis). Staging uses the TNM system.

Management depends on stage, molecular subtype, and patient factors:
\begin{itemize}[noitemsep]
  \item \textbf{Surgery:} Breast-conserving surgery (lumpectomy) with sentinel lymph node biopsy, or mastectomy (simple, modified radical, or radical). Oncoplastic techniques optimize cosmetic outcomes.
  \item \textbf{Radiation therapy:} Standard after breast-conserving surgery; post-mastectomy radiation for tumors $>$5 cm, node-positive disease, or close margins.
  \item \textbf{Systemic therapy:} Endocrine therapy (tamoxifen for premenopausal; aromatase inhibitors for postmenopausal) for ER+ disease; anti-HER2 therapy (trastuzumab, pertuzumab, T-DM1, T-DXd) for HER2+ disease; chemotherapy (anthracyclines, taxanes) for high-risk disease; CDK4/6 inhibitors for advanced ER+ disease; PARP inhibitors for BRCA-mutated tumors; immune checkpoint inhibitors for triple-negative disease with PD-L1 expression.
\end{itemize}

Prognosis varies by stage: 5-year survival is approximately 99\% for localized disease, 86\% for regional disease, and 31\% for distant metastatic disease.

\subsection{Ductal Carcinoma In Situ (DCIS)}
\label{sec:Ductal Carcinoma In Situ (DCIS)}
Ductal carcinoma in situ (DCIS) is a non-invasive neoplastic proliferation of epithelial cells confined within the mammary ductal-lobular system without penetration of the basement membrane, accounting for 20--25\% of screen-detected breast malignancies. Driven by localized genetic mutations and hormonal stimulation, DCIS is categorized by architectural patterns (comedo, cribriform, papillary, solid), with comedo DCIS presenting central high-grade necrosis and dystrophic calcification. The condition is overwhelmingly asymptomatic and detected on routine screening mammography as microcalcifications (clustered, pleomorphic, or linear branching). Diagnosis is established via stereotactic or ultrasound-guided core needle biopsy. Management involves breast-conserving surgery (lumpectomy) with negative margins ($>2$ mm) followed by radiation therapy and adjuvant tamoxifen/aromatase inhibitors for ER-positive lesions; total mastectomy is indicated for widespread or multicentric disease. Prognosis is excellent, with 20-year disease-specific survival exceeding 98\%.

\subsection{Fat Necrosis of the Breast}
\label{sec:Fat Necrosis of the Breast}
Fat necrosis of the breast is a benign, non-enzymatic inflammatory reaction resulting from focal adipocyte destruction following blunt trauma, breast surgery, core needle biopsy, or radiation therapy. Disrupted adipocytes release neutral triglycerides that undergo enzymatic saponification by tissue lipases, precipitating a foreign-body granulomatous response dominated by lipid-laden macrophages (foam cells), multinucleated giant cells, and peripheral fibrous tissue scarring. Clinical features include a solitary, firm, irregular, tethered breast mass occasionally accompanied by skin dimpling, nipple retraction, or ecchymosis, closely mimicking invasive ductal carcinoma. Diagnosis is confirmed by mammography (showing lipid cysts with eggshell calcifications or spiculated masses) and core needle biopsy when malignant features cannot be excluded. Management is conservative with patient reassurance following definitive diagnosis; symptomatic oil cysts may be aspirated. Prognosis is excellent.

\subsection{Galactocele}
\label{sec:Galactocele}
A galactocele is a benign retention cyst filled with milk or milky fluid, occurring predominantly in lactating women during or immediately following weaning. The condition results from obstruction of a lactiferous duct (secondary to inflammation, inspissated milk secretions, or external compression), leading to proximal ductal dilation and accumulation of fat-rich milk elements. Clinical features include a painless, smooth, mobile, well-circumscribed subareolar mass that may fluctuate in size with breastfeeding. Ultrasound demonstrates a well-marginated oval lesion with acoustic enhancement and a characteristic fat-fluid level. Diagnosis is established by clinical examination, sonography, and fine-needle aspiration yielding milky fluid, which is both diagnostic and therapeutic. Management involves needle aspiration for symptomatic relief; infected galactoceles require antibiotic coverage and drainage. Prognosis is excellent.

\subsection{Idiopathic Granulomatous Mastitis}
\label{sec:Idiopathic Granulomatous Mastitis}
Idiopathic granulomatous mastitis (IGM) is a rare, chronic, non-caseating inflammatory breast disease affecting parous women of childbearing age, typically within five years of lactation. The underlying pathophysiology involves an autoimmune delayed-type hypersensitivity reaction directed against extravasated proteinaceous luminal secretions from damaged ductal epithelium. Clinical features present as a rapidly enlarging, painful, hard, irregular breast mass frequently accompanied by overlying skin erythema, nipple retraction, ulceration, and chronic sterile sinus tract drainage, closely mimicking breast cancer or acute bacterial abscess. Diagnosis requires core needle biopsy demonstrating non-caseating granulomatous inflammation centered on breast lobules, with negative bacterial, fungal, and mycobacterial cultures. Management includes systemic corticosteroids (prednisone), immunosuppressive therapy (methotrexate), and conservative local care; surgical excision is reserved for refractory cases due to high recurrence rates ($20--50\%$).

\subsection{Inflammatory Breast Cancer}
\label{sec:Inflammatory Breast Cancer}
Inflammatory breast cancer is a rare, aggressive form of breast cancer (2--5\% of all breast cancers) characterized by rapid onset of breast swelling, erythema, warmth, edema (``peau d'orange''---orange peel appearance of the skin), and often a palpable mass, resulting from dermal lymphatic invasion by tumor cells. The condition is more common in younger women and African American women, with triple-negative and HER2-positive subtypes being more common than in typical breast cancer. Diagnosis is clinical, usually at an advanced stage (Stage IIIB or higher). Management is neoadjuvant chemotherapy (anthracycline- and taxane-based) followed by mastectomy and post-mastectomy radiation. Prognosis has improved but remains worse than other breast cancers, with 5-year survival approximately 40--50\%.

\subsection{Intraductal Papilloma}
\label{sec:Intraductal Papilloma}
Intraductal papilloma is a benign, focal neoplastic growth of the lactiferous ductal epithelium lining a fibrovascular core, representing the most common cause of spontaneous pathological nipple discharge. Solitary intraductal papillomas arise within major subareolar lactiferous ducts in women aged 30--50 years, whereas multiple papillomas originate in peripheral terminal duct lobular units and confer an increased risk of subsequent breast carcinoma. Torsion of the delicate fibrovascular stalk causes localized ischemia, hemorrhage, and epithelial sloughing. Clinical features present as unilateral, spontaneous, bloody (sanguineous) or serous single-duct nipple discharge, occasionally accompanied by a small subareolar palpable nodule. Diagnosis is established by galactography (ductography), ultrasound, and core needle or microdochectomy biopsy. Management involves complete surgical excision of the involved duct (microdochectomy) to rule out concurrent ductal carcinoma in situ. Prognosis is excellent.

\subsection{Invasive Lobular Carcinoma (ILC)}
\label{sec:Invasive Lobular Carcinoma (ILC)}
Invasive lobular carcinoma (ILC) is the second most common histological subtype of invasive breast cancer, accounting for 10--15\% of all cases. Pathophysiologically, ILC is characterized by loss of cellular adhesion mediated by somatic mutation or promoter hypermethylation of the E-cadherin gene (CDH1). Histologically, dyscohesive tumor cells infiltrate the stroma in a single-file, linear array ('Indian file' pattern) surrounding non-neoplastic ducts without eliciting a marked desmoplastic stroma. Consequently, ILC frequently presents as an ill-defined thickening or subtle induration rather than a discrete palpable mass, leading to delayed mammographic detection and higher rates of multifocality, bilaterality, and advanced stage at presentation. Metastatic dissemination exhibits a unique tropism for the peritoneum, retroperitoneum, gastrointestinal tract, leptomeninges, and ovaries. Diagnosis is by core biopsy. Management includes surgical resection, radiation, and long-term endocrine therapy, as most ILCs are strongly ER/PR-positive. Prognosis is comparable to invasive ductal carcinoma, though late recurrences are more frequent.

\subsection{Lobular Carcinoma In Situ (LCIS)}
\label{sec:Lobular Carcinoma In Situ (LCIS)}
Lobular carcinoma in situ (LCIS) is a non-invasive neoplastic proliferation of small, dyscohesive atypical cells completely filling and expanding the lobules and terminal ductules of the breast. Like invasive lobular carcinoma, LCIS is driven by loss of E-cadherin expression (CDH1 gene inactivation). Unlike DCIS, LCIS does not produce clinical masses or mammographic calcifications and is typically discovered incidentally on core needle biopsy performed for unrelated benign lesions. LCIS functions primarily as a powerful non-palpable, non-calcified biomarker indicating a 9- to 10-fold increased risk of developing invasive carcinoma in either breast (bilateral risk). Diagnosis requires histological confirmation with E-cadherin immunohistochemical staining. Management options include active surveillance with annual screening mammography and MRI, selective estrogen receptor modulators (tamoxifen, raloxifene) for chemoprevention, or bilateral prophylactic mastectomy for high-risk BRCA or family history carriers.

\subsection{Male Breast Cancer}
\label{sec:Male Breast Cancer}
Male breast cancer accounts for less than 1\% of all breast cancers but is increasing in incidence. Risk factors include BRCA2 mutations (most common genetic predisposition), Klinefelter syndrome, radiation exposure, gynecomastia, and liver disease, with a median age at diagnosis of 67 years. The condition results from malignant transformation of breast epithelium in males. Clinical features typically include a painless, subareolar mass, with ER-positive disease being more common than in female breast cancer. Management follows the same principles as female breast cancer: mastectomy with sentinel lymph node biopsy, adjuvant radiation, and endocrine therapy (tamoxifen is first-line for male breast cancer). Prognosis depends on stage at diagnosis.

\subsection{Mammary Duct Ectasia}
\label{sec:Mammary Duct Ectasia}
Mammary duct ectasia (periductal mastitis) is a benign inflammatory condition characterized by progressive dilation, shortening, and periductal fibrosis of the subareolar lactiferous ducts, predominantly affecting multiparous postmenopausal women. The condition results from involutional accumulation and lipid oxidation of ductal secretions, producing ductal rupture, lipid extravasation, and chronic periductal granulomatous inflammation. Clinical features include thick, sticky, pasty green or brownish single- or multiple-duct nipple discharge, subareolar pain, nipple inversion/retraction, and a firm subareolar mass mimicking breast carcinoma. Ultrasound demonstrates dilated, fluid-filled subareolar ducts containing particulate debris. Diagnosis is established by clinical examination, sonography, and biopsy to exclude malignancy. Management is conservative with warm compresses and antibiotics if infected; major duct excision (Hadfield operation) is indicated for persistent, troublesome discharge or severe pain. Prognosis is excellent.

\subsection{Paget's Disease of the Nipple}
\label{sec:Paget's Disease of the Nipple}
Paget's disease of the nipple is a rare presentation of breast cancer (1--3\% of cases) involving the nipple-areolar complex. The condition is characterized by an eczematous or psoriasiform eruption of the nipple, often with pruritus, burning, and nipple discharge, with nearly all cases being associated with an underlying breast carcinoma (in situ or invasive). Histologically, Paget cells (large, pale, intraepidermal adenocarcinoma cells) are found within the epidermis. Diagnosis requires biopsy of the affected nipple. Management typically involves mastectomy with sentinel lymph node biopsy, or breast-conserving surgery if the underlying carcinoma can be adequately excised. Prognosis depends on the stage of the underlying carcinoma rather than the Paget's disease itself.

\subsection{Triple-Negative Breast Cancer (TNBC)}
\label{sec:Triple-Negative Breast Cancer (TNBC)}
Triple-negative breast cancer (TNBC) is a highly aggressive molecular subtype of breast cancer defined by the complete lack of estrogen receptor (ER), progesterone receptor (PR), and human epidermal growth factor receptor 2 (HER2) expression, accounting for 10--20\% of all breast malignancies. TNBC occurs with higher frequency in younger women, women of African descent, and individuals harboring germline BRCA1 mutations. Lacking targeted hormonal or HER2 pathways, TNBC exhibits rapid cellular proliferation (high Ki-67 index), high histological grade, and early hematogenous dissemination, with a peak incidence of visceral (lung, liver) and central nervous system metastases within 3 years of diagnosis. Diagnosis is confirmed by immunohistochemistry ($<1\%$ ER/PR staining) and HER2 FISH testing. Management involves neoadjuvant combination chemotherapy (anthracyclines, taxanes, platinum agents) combined with immune checkpoint inhibitors (pembrolizumab) for PD-L1-positive tumors, followed by definitive surgery and post-neoadjuvant capecitabine or PARP inhibitors (olaparib) for BRCA carriers. Prognosis is worse than receptor-positive subtypes, with a 5-year survival rate of approximately 77\%.

